\begin{explanationbox}
	\textbf{\textcolor{CExplanation}{一句话直觉}}

	二项式展开的本质是“从 $n$ 个 $(a+b)$ 中选出 $r$ 个提供 $b$，其余提供 $a$”。

	\medskip
	\textbf{\textcolor{CExplanation}{核心拆解}}
	\begin{description}[style=nextline, leftmargin=2.8em, labelsep=0.5em, itemsep=0.45em]
		\item[\textbf{要点一（组合选择方式）：}] 每一项中的 $b$ 来自 $n$ 个因式中的某 $r$ 个，每个因式要么选 $a$ 要么选 $b$，且选定后 $a$ 的指数为 $n-r$，$b$ 的指数为 $r$。
		\item[\textbf{要点二（项序与 $r$ 对应）：}] 公式中 $T_{r+1}$ 指出，$r$ 决定 $b$ 的指数，同时决定它是展开式的第 $r+1$ 项，$r$ 从 $0$ 开始。
	\end{description}

	\medskip
	\textbf{\textcolor{CExplanation}{几何本质（离散分配结构）}}

	展开式每一项对应一种从 $n$ 个位置中选择 $r$ 个分配 $b$ 的方式，组合数 $\binom{n}{r}$ 正是这种分配方式的个数。

	\medskip
	\textbf{\textcolor{CExplanation}{代数意义（乘积展开）}}

	$(a+b)^n$ 的展开是 $n$ 个一次二项式连乘的结果，通项公式将展开后的同类项合并系数用组合数表示，实现从因式积到多项式和的转化。

	\medskip
	\textbf{\textcolor{CExplanation}{考点价值（定向求项）}}
	\begin{itemize}[leftmargin=*, itemsep=0.5em, topsep=0.4em]
		\item \textbf{考法一（指定项求值）：} 直接代入 $r$ 求特定项的系数或项本身，如常数项、含 $x^k$ 的项。
		\item \textbf{考法二（系数性质判断）：} 利用二项式系数的对称性、单调性（在中间项取得最大值）进行推理。
	\end{itemize}

	\medskip
	\textbf{\textcolor{CExplanation}{顿悟点}}

	\textbf{通项公式是“谁展开、第几项、长什么样”三问的合一：给定 $n$、$r$ 即唯一确定一项。}

	\medskip
	\textbf{\textcolor{CExplanation}{使用场景}}
	\begin{itemize}[leftmargin=*, itemsep=0.5em, topsep=0.4em]
		\item \textbf{场景一（求常数项）：} 先写出并化简通项，再令整个通项中变量的指数为 $0$，由此确定 $r$。
		\item \textbf{场景二（求有理项）：} 要求展开式中所有指数为整数的项，通过 $r$ 满足某种整除条件。
	\end{itemize}
\end{explanationbox}